
\documentclass{article} %210 mm × 297 mm
\usepackage[utf8]{inputenc}
\usepackage[a4paper,left=1.5cm, right=1.5cm, top=2cm, bottom=2cm]{geometry}
\usepackage{graphicx}
\usepackage{systeme}          % for Solution 3
\usepackage{array}  
%\usepackage{blindtext}
\usepackage{multirow}
\usepackage{mhchem}
\usepackage{url}
\usepackage{subfigure}
\usepackage{amsfonts,latexsym} % para tener disponibilidad de diversos simbolos
\usepackage{enumerate}
%\usepackage{booktabs}
\usepackage{float}
%\usepackage{threeparttable}
\usepackage{array,colortbl}
%\usepackage{ifpdf}
%\usepackage{rotating}
\usepackage{stfloats}
\usepackage{url}
\usepackage{listings}
\usepackage{fancyhdr}
\usepackage{biblatex}
\usepackage{color}
\addbibresource{sources.bib}
%\usepackage{hyperref}
%\usepackage{wrapfig}
\usepackage{array}
%\usepackage{multirow}
%\usepackage{adjustbox}
%\usepackage{nccmath}
\usepackage[dvipsnames]{xcolor}
\usepackage{tcolorbox}
\newtcolorbox{mybox}{colback=white,
colframe=black}
\newcommand{\titulo}{Abgabe 9; Diskrete Mathematik}
\newcommand{\nombre}{Liliana Sanfilippo}
\newcommand{\FCE}{\textbf{WiSe 2024/2025}}

 

\usepackage{fancyhdr}
\pagestyle{fancy}
\fancyhead{}
\fancyfoot{}
\fancyhead[LO]{\small\nombre}
\fancyhead[RE]{\small\titulo}
\fancyhead[LO]{\small\FCE}
\fancyfoot[RO,LE] {\thepage}
\usepackage[onehalfspacing]{setspace}

\setcounter{page}{1} %Número de la página, la editorial lo cambiará
\begin{document}
\begin{tabular}{p{12cm}}
{{\Large\bf\titulo}}\\
\vspace{0.2cm}\\
 {\large\nombre}\\
 \vspace{0.1cm}
\end{tabular}
\section*{Präsenzaufgaben}

\textbf{1. Welche der folgenden Codes sind linear? Bestimmen Sie jeweils den Minimalabstand und wie viele Fehler erkannt bzw. korrigiert werden können.}
\begin{itemize}
  \item[(a)] \( \{0000, 1100, 1010, 1001, 0110, 0101, 0011, 1111\} \)
  \begin{mybox}
  \paragraph{Nullvektor:}
  Ist vorhanden (0000) 
  \paragraph{Abgeschlossenheit unter Addition:}
Mit dem Nullvektor addiert bleiben alle Codes sie selbst, also nicht in der Tabelle.  
  \begin{table}[H]
      \centering
      \begin{tabular}{c|ccccccc}
        & 1100 & 1010 & 1001 & 0110 & 0101 & 0011 & 1111\\ \hline
    1100& 0000 & 0110 & 0101 & 1010 & 1001 & 1111 & 0011\\
    1010&  & 0000 & 0110 & 1100 & 1111 & 1001 & 0101\\
    1001&  &  & 0000 & 1111 & 1100 & 1010 & 0110\\
    0110&  &  &  & 0000 & 1010 & 0101 & 1001\\
    0101&  &  &  &  & 0000 & 0110 & 1010 \\
    0011&  &  &  &  &  & 0000 & 1100\\
    1111&  &  &  &  &  &  & 0000
      \end{tabular}
  \end{table}
  Der Code ist also linear. 
  \paragraph{Minimalabstand:} 2
  \paragraph{t-Fehlererkennend:} für $2 > t$ (bis zu 1 Fehler) 
  \paragraph{t-Fehlerkorrigierend:} $2 > 2t$ (keine Fehler)
  \end{mybox}
  \item[(b)] \( \{10000, 01010, 00001\} \)
\begin{mybox}
Ist nicht linear, da er keinen Nullvektor beinhaltet. 
\paragraph{Minimalabstand:} 2
\paragraph{t-Fehlererkennend:} für $2 > t$ (bis zu 1 Fehler) 
\paragraph{t-Fehlerkorrigierend:} für $2 > 2t$ (keine Fehler) 
\end{mybox}
  \item[(c)] \( \{00000000, 10101010, 01010101\} \)
  \begin{mybox}
  Ist nicht linear, da nicht abgeschlossen unter Addition: 
  \[10101010 + 01010101 = 11111111 \notin \{00000000, 10101010, 01010101\} \]
  \paragraph{Minimalabstadn:} 4
  \paragraph{t-Fehlererkennend:} für $4 > t$ (bis zu 3 Fehler) 
  \paragraph{t-Fehlerkorrigierend:} $4 > 2t$ (bis zu 1 Fehler) 
  \end{mybox}
\end{itemize}



\section*{Aufgaben zur Abgabe}

\subsection*{Aufgabe 1 (4 Punkte)}
Sei $S$ eine Differenzenfamilie in $\mathbb{Z}/m\mathbb{Z}$. Zeigen Sie, dass $S + i$ eine Differenzenfamilie in $\mathbb{Z}/m\mathbb{Z}$ ist für jedes $i \in \mathbb{Z}/m\mathbb{Z}$. Folgern Sie, dass man ohne Einschränkung annehmen kann, dass jede Differenzenfamilie $0$ und $1$ enthält.
\begin{mybox}
Sei $S + i = \{s+i \mod m \mid s \in S\}$, also eine Verschiebung von $S$. Da für $s_1 \neq s_2; s_1, s_2 \in S$ gilt für $S+i$: $(s_1 + i) - (s_2 + i) \mod m = s_1 - s_2 \mod m$. \\
Weil $S$ eine Differenzenfamilie modulo $m$ ist und jedes $i$ auch aus einer kommt, ist auch $S + i$ eine Differenzenfamilie modulo $m$ und somit in $\mathbb{Z}/m\mathbb{Z}$. \\ \newline 
Für alle $S =\{ s_1, s_2, ...\}$ kann man eine Verschiebung der Art $-s_1$ erzeugen, so dass $S' = \{0, s_2 - s_1, ...\}$ bzw. $S' = \{0, s_2', ...\}$. Also liegt 0 in jeder Differenzenfamilie. \\
Wählt man ein $x \in \mathbb{Z}/m\mathbb{Z}$ mit $ggT(x,m) = 1$ (invertierbar), kann man jedes $s_i \in S'$ mit $x$ multiplizieren und erhält $S'' = \{0 \mod m, s_2'\cdot x \mod m, ... \} = \{0, 1, ...\}$.  
\end{mybox}

\subsection*{Aufgabe 2 (4 Punkte)}
Sei $S$ eine Differenzenfamilie in $\mathbb{Z}/m\mathbb{Z}$. Zeigen Sie, dass $T = \mathbb{Z}/m\mathbb{Z} \setminus S$ eine Differenzenfamilie ist. Bestimmen Sie die Parameter des von $T$ induzierten $2$-Designs in Abhängigkeit der Parameter des von $S$ induzierten $2$-Designs.
\begin{mybox}
Da
\[\# \text{Anzahl Restklassen in } \mathbb{Z}/m\mathbb{Z} = m \]
gilt
\[|T| = m - |S|\]
Ein $x \in \mathbb{Z}/m\mathbb{Z} \backslash \{0\}$ trifft für $S$ genau $a_S$ mal und für $T$ genau $a_T$ mal. 
Da 
\[\# \text{Anzahl Restklassen in } \mathbb{Z}/m\mathbb{Z} \backslash \{0\} = m -1 = \# \text{Gesamtanzahl die $x$ trifft}\]
Also 
\begin{align*}
    a_S + a_T & = m -1 \\
     a_T & = m -1 - a_S 
\end{align*}
Also ist $T$ eine Differenzenfamilie, bei der jede Differenz $a_T$-mal vorkommt. 
\\
Für ein 2-Design sind die Parameter bein Induzierung durch $S$: 
\[v = |S|, \quad k = |S|, \quad \lambda = a_S\]
Aufgrund der Komplementarität von $S$ und $T$  sind die Parameter für ein durch $T$ induziertes 2-Design: 
\[v = |T| = m - |S|, \quad k = |T| = m - |S|, \quad \lambda = a_T =  m-1-a_S\]
\end{mybox}

\subsection*{Aufgabe 3 (4 Punkte)}
Ergänzen Sie die Menge von Codewörtern
\[
\{00000, \ 00110, \ 10101\}
\]
zu einem linearen Code $C$. Finden Sie eine Kontrollmatrix für $C$ und bestimmen Sie das Gewicht aller Codewörter in $C$.
\begin{mybox}
\begin{table}[H]
    \centering
    \begin{tabular}{c|cccc}
        & 00000 & 00110 & 10101 & 10011\\ \hline
   00000& 00000 & 00110 & 10101 & 10011\\
   00110&       & 00000 & 10011 & 10101\\
   10101&       &       & 00000 & 00110\\
   10011&       &       &       & 00000\\
    \end{tabular}
\end{table}
Ergänzter Code: \{00000, \ 00110, \ 10101, \ 10011\}\\
Gewichte: 
\begin{align*}
    00000: & \quad 0\\
    00110: & \quad 2\\
    10101: & \quad 3\\
    10011: & \quad 3
\end{align*}
\end{mybox}
\begin{mybox}
\[
Generatormatrix = \begin{pmatrix}
0 & 0 & 0 & 0 & 0 \\
0 & 0 & 1 & 1 & 0 \\
1 & 0 & 1 & 0 & 1 \\
1 & 1 & 1 & 0 & 0
\end{pmatrix}
\]
Wie man die Kontrollmatrix berechnet habe ich nicht verstanden. 
\end{mybox}
\subsection*{Aufgabe 4 (4 Punkte)}
Der Code $C$ sei durch die folgende Kontrollmatrix gegeben:
\[
\begin{pmatrix}
1 & 0 & 1 & 0 & 1 & 1 \\
0 & 1 & 0 & 0 & 1 & 1 \\
1 & 0 & 0 & 1 & 1 & 0 \\
\end{pmatrix}
\]
\begin{enumerate}
    \item[(a)] Korrigieren Sie folgende empfangenen Wörter zu Codewörtern in $C$ unter der Annahme, dass höchstens ein Fehler aufgetreten ist:
    \[
    110111, \quad 101010, \quad 110101.
    \]
    \begin{mybox}
    \[
    110111: 
\begin{pmatrix} 
1 & 0 & 1 & 0 & 1 & 1 \\
0 & 1 & 0 & 0 & 1 & 1 \\
1 & 0 & 0 & 1 & 1 & 0
\end{pmatrix} 
\cdot 
\begin{pmatrix} 
1 \\
1 \\
0 \\
1 \\
1 \\
1
\end{pmatrix}
= \begin{pmatrix}
1 \cdot 1 + 0 \cdot 1 + 1 \cdot 0 + 0 \cdot 1 + 1 \cdot 1 + 1 \cdot 1 \\
0 \cdot 1 + 1 \cdot 1 + 0 \cdot 0 + 0 \cdot 1 + 1 \cdot 1 + 1 \cdot 1 \\
1 \cdot 1 + 0 \cdot 1 + 0 \cdot 0 + 1 \cdot 1 + 1 \cdot 1 + 0 \cdot 1
\end{pmatrix}
= \begin{pmatrix}
1 \\
1 \\
1
\end{pmatrix}
\]
$ 110111 \rightarrow 110101$
\[
101010: 
\begin{pmatrix} 
1 & 0 & 1 & 0 & 1 & 1 \\
0 & 1 & 0 & 0 & 1 & 1 \\
1 & 0 & 0 & 1 & 1 & 0
\end{pmatrix}
\cdot 
\begin{pmatrix} 
1 \\
0 \\
1 \\
0 \\
1 \\
0
\end{pmatrix}
= 
\begin{pmatrix} 
1 \cdot 1 + 0 \cdot 0 + 1 \cdot 1 + 0 \cdot 0 + 1 \cdot 1 + 1 \cdot 0 \\
0 \cdot 1 + 1 \cdot 0 + 0 \cdot 1 + 0 \cdot 0 + 1 \cdot 1 + 1 \cdot 0 \\
1 \cdot 1 + 0 \cdot 0 + 0 \cdot 1 + 1 \cdot 0 + 1 \cdot 1 + 0 \cdot 0
\end{pmatrix}
= 
\begin{pmatrix} 
1 \\
1 \\
0
\end{pmatrix}
\]
$101010 \rightarrow 101011$
\[
110101: 
\begin{pmatrix} 
1 & 0 & 1 & 0 & 1 & 1 \\
0 & 1 & 0 & 0 & 1 & 1 \\
1 & 0 & 0 & 1 & 1 & 0
\end{pmatrix}
\cdot 
\begin{pmatrix} 
1 \\
1 \\
0 \\
1 \\
0 \\
1
\end{pmatrix}
= 
\begin{pmatrix} 
1 \cdot 1 + 0 \cdot 1 + 1 \cdot 0 + 0 \cdot 1 + 1 \cdot 0 + 1 \cdot 1 \\
0 \cdot 1 + 1 \cdot 1 + 0 \cdot 0 + 0 \cdot 1 + 1 \cdot 0 + 1 \cdot 1 \\
1 \cdot 1 + 0 \cdot 1 + 0 \cdot 0 + 1 \cdot 1 + 1 \cdot 0 + 0 \cdot 1
\end{pmatrix}
= 
\begin{pmatrix} 
0 \\
0 \\
0
\end{pmatrix}
\]
Bleibt so, ist richtig. 
    \end{mybox}
    \item[(b)] Wie viele Wörter existieren, die nicht durch Änderung von höchstens einem Bit zu einem Codewort in $C$ korrigiert werden können?
\end{enumerate}
\begin{mybox}
Dafür erstmal den Code berechnen: 
\[
\begin{aligned}
1x_1 + 0x_2 + 1x_3 + 0x_4 + 1x_5 + 1x_6 &= 0 \\
0x_1 + 1x_2 + 0x_3 + 0x_4 + 1x_5 + 1x_6 &= 0 \\
1x_1 + 0x_2 + 0x_3 + 1x_4 + 1x_5 + 0x_6 &= 0
\end{aligned}
\]
\[
\left( \begin{array}{c cccccc|c}
I. &1 & 0 & 1 & 0 & 1 & 1 & 0 \\
II. & 0 & 1 & 0 & 0 & 1 & 1 & 0 \\
III. & 1 & 0 & 0 & 1 & 1 & 0 & 0
\end{array} \right)
\]
III. - I. 
\[
\left( \begin{array}{c cccccc|c}
I. & 1 & 0 & 1 & 0 & 1 & 1 & 0 \\
II. & 0 & 1 & 0 & 0 & 1 & 1 & 0 \\
III. & 0 & 0 & 1 & 1 & 0 & 1 & 0
\end{array} \right)
\]
I. - III. 
\[
\left( \begin{array}{c cccccc|c}
I.      & 1 & 0 & 0 & 1 & 1 & 0 & 0 \\
II.     & 0 & 1 & 0 & 0 & 1 & 1 & 0 \\
III.    & 0 & 0 & 1 & 1 & 0 & 1 & 0
\end{array} \right)
\]
II. - III. 
\[
\left( \begin{array}{c cccccc|c}
I.      & 1 & 0 & 0 & 1 & 1 & 0 & 0 \\
II.     & 0 & 1 & 1 & 1 & 1 & 1 & 0 \\
III.    & 0 & 0 & 1 & 1 & 0 & 1 & 0
\end{array} \right)
\]
\[
\left( \begin{array}{c cccccc|c}
I.      & 1 & 0 & 0 & 1 & 1 & 0 & 0 \\
II.     & 0 & 1 & 1 & 1 & 1 & 1 & 0 \\
III.    & 0 & 0 & 1 & 1 & 0 & 1 & 0
\end{array} \right)
\]
II. 
\begin{align*}
    x_2+x_5+x_6& =0 \\
    x_5+x_6& = x_2
\end{align*}
III. 
\begin{align*}
    x_1+x_4+x_5& =0 \\
    x_4+x_5& = x_1
\end{align*}
I.
\begin{align*}
   x_1+x_3+x_5+x_6 & =0 \\
   (x_4+x_5)+x_3+x_5+x_6 & =0 \\
   x_4+x_3+2x_5+x_6 & =0 \\
   x_4+x_3+x_6 & =0 \\
   x_4+x_6 & = x_3
\end{align*}
Codewörter: 
\[
(x_1, x_2, x_3, x_4, x_5, x_6) = (x_4 + x_5, x_5 + x_6, x_4 + x_6, x_4, x_5, x_6)
\]
$\{(0,0,0,0,0,0), (0,0,1,0,0,1), (1,1,0,0,1,0), (1,0,1,0,1,1), (1,0,1,1,0,0), (1,1,0,1,0,1)
(0,1,1,1,1,0), (0,0,0,1,1,1)\}$
\end{mybox}
\begin{mybox}
Der Minimalabstand des Codes ist 3 (siehe bspw. (1,0,1,0,1,1)  und (0,1,0,0,1,1)). Da es $2^6 = 64$ mögliche 6-stellige Wörter aus 0,1 gibt und Wörter mit einem Bit nicht korrigiert werden können (davon gibt es pro Wort 6: $6 \cdot 8$) und der Code 8 Wörter hat: 
\[64-8-6\cdot 8=48\]
48 Wörter können nicht mit nur einem Bit korrigiert werden. 
\end{mybox}
\end{document}